% 1.2.3
\subsubsection{Đánh giá ưu nhược điểm của RSA và AES}
\indent Hai thuật toán RSA và AES là những đại diện tiêu biểu cho hai hướng tiếp cận khác nhau trong lĩnh vực mật mã học. Mặc dù đều nhằm mục tiêu bảo vệ thông tin khỏi truy cập trái phép, song chúng khác biệt rõ rệt về cấu trúc, cơ chế hoạt động, cũng như hiệu năng xử lý.

\indent Thuật toán RSA được xây dựng trên cơ sở toán học của phép nhân và phân tích thừa số nguyên tố lớn, sử dụng hai khóa khác nhau: khóa công khai (public key) để mã hóa và khóa bí mật (private key) để giải mã. Ưu điểm nổi bật của RSA là độ an toàn cao, khó bị tấn công bằng phương pháp phân tích ngược, và giải quyết tốt vấn đề xác thực và đảm bảo nguồn gốc dữ liệu. RSA thường được ứng dụng trong các hệ thống yêu cầu xác minh danh tính hoặc ký số, nơi tính bảo mật được ưu tiên hàng đầu. Tuy nhiên, hạn chế của RSA nằm ở hiệu suất xử lý chậm, đặc biệt khi kích thước khóa lớn, và khả năng mở rộng kém khi mã hóa dữ liệu dung lượng lớn. Việc sinh khóa và thực hiện phép mũ modulo trên các số lớn cũng làm tiêu tốn đáng kể tài nguyên tính toán.

\indent Trong khi đó, AES thuộc nhóm mã hóa đối xứng, sử dụng cùng một khóa duy nhất cho cả hai quá trình mã hóa và giải mã. AES hoạt động dựa trên các vòng lặp hoán vị và thay thế (Substitution–Permutation Network) giúp dữ liệu sau khi mã hóa trở nên khó phân tích và gần như không thể khôi phục nếu không có khóa hợp lệ. Ưu điểm chính của AES là tốc độ xử lý cao, hiệu năng ổn định, và tính linh hoạt khi triển khai trên cả phần cứng lẫn phần mềm. Nhờ đó, AES được sử dụng phổ biến trong nhiều hệ thống cần mã hóa dữ liệu theo thời gian thực như truyền thông mạng, mã hóa ổ đĩa, hoặc bảo vệ cơ sở dữ liệu. Dù vậy, điểm yếu lớn nhất của AES là rủi ro khi quản lý khóa: vì cùng một khóa dùng cho cả hai phía, nên nếu khóa bị lộ, toàn bộ dữ liệu có thể bị giải mã.

\indent Nhìn chung, RSA mạnh về bảo mật và xác thực, trong khi AES vượt trội về tốc độ và hiệu quả xử lý. RSA thích hợp cho các tình huống cần đảm bảo an toàn tuyệt đối trong trao đổi thông tin nhạy cảm, còn AES phù hợp với các ứng dụng yêu cầu hiệu năng cao và mã hóa khối lượng dữ liệu lớn. Việc lựa chọn thuật toán nào phụ thuộc vào mục tiêu bảo mật và mức độ ưu tiên giữa tốc độ và độ an toàn trong từng ứng dụng cụ thể.
